% Pass: Opus 5 (claude-opus-5), 2026-09-06 - first pass, unchecked against the
% sheets by a human. Its own context, no batch before it in this conversation.
%
% Book I, batch 7: sheets 73 to 84, leaves 64 to 71. Leaf 66 is photographed
% twice and two sheets carry no leaf number at all - a loose photograph with
% its verso, and a separate sheet of thumbnail drawings laid over leaf 69 - so
% the batch covers eight leaves and three unnumbered sheets.
%
% Five things worth knowing before reading it.
%
% The subject does not change once in this batch, which is unusual for Book I.
% Leaves 64 to 71 are eight consecutive leaves of watercolours, running from
% the Rockland and Gloucester summer of 1926 to the Cape Elizabeth and Essex
% summer of 1929, and every leaf is ruled the same way: a delivery date, a
% price, a description, the amount received net of the third that went to
% Rehn, and the date the cheque cleared. Almost every picture is 14 x 20 and
% almost every price is 250, 300 or 400.
%
% The headings run the way batch 6 found them and not the way batch 5 did.
% Leaves 64, 65 and 66 head their money columns in ink and their delivery
% column in pencil - « Delivered to Rehn », « Price », « Amt. rec'd. »,
% « When rec'd. »; leaf 68 heads one column only, « When payment rec'd. »; and
% leaves 67, 69, 70 and 71 head nothing. No headings are supplied in \add{}
% anywhere in this batch.
%
% Leaf 66 is the hinged-clipping case again and it costs nothing this time. On
% sheet 17992 a large halftone of the Custom House, Portland lies over the
% lower half of the leaf; on sheet 16820, which the Whitney calls « Page 66 -
% Verso », it has been turned back clear of the writing altogether. The
% transcription is at 16820, where the leaf is whole, and the clipping and a
% note stand at 17992.
%
% Two sheets are not leaves and are not versos of leaves. Sheets 18051 and
% 17242 are the two faces of a loose photograph of the Manhattan Bridge
% Entrance watercolour, laid between leaves 64 and 65; its verso carries four
% lines naming the owner, and they agree with the last entry of leaf 65, which
% is the same picture. Sheet 18337 is a separate sheet of eleven ink thumbnail
% drawings with their sale lines, lying over leaf 69 - the Whitney's descriptor
% says it is adhered there - and it is a document in its own right, not a
% second photograph of the leaf underneath. Both sheets are transcribed.
%
% Leaf 71 is the one place this batch cannot finish. The Farm at Essex
% halftone is hinged over the lower left of the leaf, covering two entire
% blocks, and the Whitney's second photograph of that leaf - sheet 16955, the
% recto again with the clipping turned back - is sheet 85, the first sheet of
% batch 8. So leaf 71 is transcribed here as far as this batch's photograph
% shows it, the covered rows are \ill{}, and the note says where they will be
% read. Batch 8 should add those two blocks at 16955 and must not transcribe
% the four blocks already given here.
\documentclass[11pt,a4paper]{article}
% The preamble shared by every transcription of the Hopper artist's ledgers.
%
% It rests on one idea, taken whole from the Grothendieck workbench this
% method comes from: **uncertainty compounds, it does not resolve.** A
% transcription that smooths over a doubtful figure has destroyed the only
% thing it was made for — a reader must be able to tell, at every point, what
% was on the leaf and what was supplied.
%
% A ledger sharpens that. In a manuscript of mathematics an invented word is a
% wrong word. Here an invented figure is a *sale that did not happen*, at a
% price nobody paid, to a buyer who never bought — and it will be cited,
% because a table looks like data in a way that prose does not.
%
% The same commands are recognised by `scripts/render.mjs`, which produces the
% site's reading view. A macro added here must be added there too, or rendering
% fails loudly, which is the wanted behaviour: a silently wrong rendering is
% worse than none.

% \ProvidesFile, not \ProvidesPackage: the transcriptions pull this in with
% % The preamble shared by every transcription of the Hopper artist's ledgers.
%
% It rests on one idea, taken whole from the Grothendieck workbench this
% method comes from: **uncertainty compounds, it does not resolve.** A
% transcription that smooths over a doubtful figure has destroyed the only
% thing it was made for — a reader must be able to tell, at every point, what
% was on the leaf and what was supplied.
%
% A ledger sharpens that. In a manuscript of mathematics an invented word is a
% wrong word. Here an invented figure is a *sale that did not happen*, at a
% price nobody paid, to a buyer who never bought — and it will be cited,
% because a table looks like data in a way that prose does not.
%
% The same commands are recognised by `scripts/render.mjs`, which produces the
% site's reading view. A macro added here must be added there too, or rendering
% fails loudly, which is the wanted behaviour: a silently wrong rendering is
% worse than none.

% \ProvidesFile, not \ProvidesPackage: the transcriptions pull this in with
% % The preamble shared by every transcription of the Hopper artist's ledgers.
%
% It rests on one idea, taken whole from the Grothendieck workbench this
% method comes from: **uncertainty compounds, it does not resolve.** A
% transcription that smooths over a doubtful figure has destroyed the only
% thing it was made for — a reader must be able to tell, at every point, what
% was on the leaf and what was supplied.
%
% A ledger sharpens that. In a manuscript of mathematics an invented word is a
% wrong word. Here an invented figure is a *sale that did not happen*, at a
% price nobody paid, to a buyer who never bought — and it will be cited,
% because a table looks like data in a way that prose does not.
%
% The same commands are recognised by `scripts/render.mjs`, which produces the
% site's reading view. A macro added here must be added there too, or rendering
% fails loudly, which is the wanted behaviour: a silently wrong rendering is
% worse than none.

% \ProvidesFile, not \ProvidesPackage: the transcriptions pull this in with
% \input{../preamble/hopper} rather than \usepackage, and \ProvidesPackage in
% an \input-ed file makes LaTeX warn on every compile that it "requested
% package `'" and got `hopper'. A warning nobody can act on is a warning
% everybody learns to scroll past, which is how the real ones get missed.
\ProvidesFile{hopper.sty}[2026/09/05 Transcriptions of the Hopper artist's ledgers]

\usepackage{array}
% longtable, not tabularx: a leaf of thirty-one exhibition lines must break
% across pages, and tabularx cannot be wrapped in an environment of our own —
% it scans the source for a literal \end{tabularx} and dies on \end{ledgertable}
% with « File ended while scanning use of \TX@get@body ». longtable has no such
% scan, and the wrapping column type below does the job tabularx's X would.
\usepackage{longtable}

% A wrapping column, `Y{0.3}` being three tenths of the measure.
%
% Prose columns must be declared this way and not as `l`. A table of `l`
% columns is set to its natural width, which can be wider than the page, and
% TeX issues **no warning at all** because nothing ever asked the row to fit:
% the first compile of Book I produced a page whose right-hand column ran off
% the paper with a clean log. A fixed-width column wraps, and if it still
% cannot fit, TeX says so — which is what `npm run pdf` then refuses to build
% over.
%
% Leave about 0.03 of the measure per column for the inter-column space: five
% columns can share 0.85, six can share 0.80. Getting it wrong is not a
% guessing game — `npm run pdf` reports the overfull box and how many points
% too wide it is.
\newcolumntype{Y}[1]{>{\raggedright\arraybackslash}p{#1\linewidth}}
\usepackage[english]{babel}
\usepackage{fontspec}
\usepackage{geometry}
\geometry{a4paper,margin=25mm,marginparwidth=28mm}
\usepackage{marginnote}
\usepackage{xcolor}
\usepackage{tikz}
% ulem, not soul. `soul`'s \st and \ul are built on a character-by-character
% scan that XeLaTeX's font machinery breaks: the document fails with an
% undefined control sequence rather than degrading. `normalem` stops it
% hijacking \emph, which the transcriptions use for underlining of Jo Hopper's
% own.
\usepackage[normalem]{ulem}
% ulem sets each underlined word in a box TeX can neither hyphenate nor break
% after, so a paragraph dense in \uncertain{} readings can reach a state where
% no legal breakpoint exists. \emergencystretch lets TeX loosen it on a third
% pass instead of overrunning: an overfull line in a transcription hides
% content.
\setlength{\emergencystretch}{3em}
\usepackage{hyperref}
\hypersetup{colorlinks=true,linkcolor=black,urlcolor=black,pdfborder={0 0 0}}

% --- Batch metadata ---------------------------------------------------------
% Taken as they stand by the rendering script, which shows them at the head of
% the reading view. They fix what the file claims to cover: without them a
% transcription floats unmoored in an archive of five hundred sheets.

\newcommand{\ledger}[1]{\gdef\hopLedger{#1}}
\newcommand{\ledgertitle}[1]{\gdef\hopLedgerTitle{#1}}
% A notebook of Josephine Hopper's, in place of a ledger: one file to a
% notebook, filed under transcripts/notebooks/, the argument the site's slug
% (« battle-of-wash-sq »). Sets the ledger to « notebooks » for everything
% that reads it, so the title block and the watermark behave as they do for a
% batch; the rendering checks the sheets against the notebook's own listing.
\newcommand{\notebook}[1]{\gdef\hopLedger{notebooks}\gdef\hopNotebook{#1}}
\gdef\hopNotebook{}
% The Whitney's accession number for the physical volume — 96.208 … 96.213.
% This is what a citation of the object cites.
\newcommand{\objectnumber}[1]{\gdef\hopObject{#1}}
% The batch this file covers. Leaving it out drops the « (batch N) » from the
% title block rather than printing a number that would be false.
\newcommand{\batch}[1]{\gdef\hopBatch{#1}}
\newcommand{\sheets}[2]{\gdef\hopFirst{#1}\gdef\hopLast{#2}}

% The Whitney's date range for the volume, copied verbatim from its resource
% record — never inferred here, never narrowed to the years this batch
% happens to cover.
\newcommand{\dating}[1]{\gdef\hopDating{#1}}

% The legal notice, printed in the title block and repeated as a diagonal
% watermark on every page of the PDF.
%
% This is not boilerplate and it is not optional. The ledgers are © Heirs of
% Josephine N. Hopper, licensed by the Artists Rights Society; the Whitney
% records the object rights as transferred to the Museum. A transcription
% reproduces Josephine Hopper's text. Every file must therefore say what it is
% on its own face, not only in a README that does not travel with the PDF.
% A file without this line gets no watermark, so omitting it is visible in the
% output rather than silent.
\newcommand{\watermark}[1]{\gdef\hopWatermark{#1}}

\gdef\hopLedger{?}\gdef\hopLedgerTitle{}\gdef\hopObject{}\gdef\hopBatch{}%
\gdef\hopFirst{?}\gdef\hopLast{?}\gdef\hopDating{}\gdef\hopWatermark{}

% --- Keywords ---------------------------------------------------------------
%
% Closing a transcription: the vocabulary under which to search for what these
% leaves record — dealers, exhibiting societies, media, places.
% `npm run manifest` extracts them to tag the whole ledger. This line is the
% single source of the tags; there is no tags file, so no tag can describe
% leaves nobody has read.
%
% It is the one thing in the file that is not on the paper, and it earns its
% place: a reader looking for Keppel needs some way in, and a tag written by
% whoever read the sheets is the only kind that cannot describe unread ones.
\newcommand{\keywords}[1]{\par\smallskip\noindent
  {\footnotesize\textsc{Keywords} --- \textit{#1}}\par}

% --- The anchor -------------------------------------------------------------

% The sheet begins here. Two arguments, and both are needed.
%
% #1 is the Whitney's **resource ref** — 16853 — which is the sheet's whole
% address: it is what the image URL is built from and what turns the facsimile
% as the transcript is scrolled. #2 is the number **written on the paper**, or
% empty where nobody wrote one, which is the case for every cover, flyleaf,
% index and loose insertion — about one sheet in nine.
%
% Keeping both is the whole of the numbering problem in this archive. Three
% sequences overlap: the Hoppers' own leaf numbers, the Whitney's order of
% photographs, and the resource refs, which are in no order at all. Only the
% first is citable and only the third is addressable, so a transcription that
% recorded one of them would be either uncitable or unopenable.
%
% `scripts/render.mjs` checks #2 against what `scripts/catalogue.mjs` parsed
% out of the Whitney's descriptor and fails on a disagreement, because a
% mistyped leaf number silently moves an entry to another year.
\newcommand{\sheet}[2]{%
  \marginnote{\footnotesize\color{gray!70}\textsf{\ifx\relax#2\relax---\else#2\fi}}%
  \label{sheet:#1}\ignorespaces}

% --- The critical apparatus -------------------------------------------------

% Illegible: never guessed. On a ledger leaf this is most often a figure, and
% a guessed figure is a fabricated transaction.
\newcommand{\ill}{\textcolor{red!60!black}{[\,\ldots\,]}}

% A doubtful reading: the word or figure is given, and flagged as uncertain.
%
% Underlining belongs to this macro and to nothing else. Jo Hopper underlines
% her own headings --- « Etchings », « Etching Plates », « Early oil » --- and
% those are set with \emph{} instead, because a reader must be able to tell a
% doubtful reading from her emphasis at a glance and the two cannot both be
% underlines. Which of them keeps the underline is decided by consequence: an
% emphasis shown as italic loses nothing, and a doubtful figure shown as a
% certain one loses everything.
\newcommand{\uncertain}[1]{\uline{#1}}

% An editorial addition — an expanded abbreviation, an implied dollar sign.
\newcommand{\add}[1]{\textcolor{blue!50!black}{[#1]}}

% Struck out in the book. Jo Hopper struck a great deal: a price revised, a
% buyer who withdrew, an exhibition that fell through. What she crossed out is
% part of the record and often the more interesting half of it.
\newcommand{\struck}[1]{\sout{#1}}

% The ink it is written in, where the ink says something.
%
% Book IV is a state machine and the colour is its status field: charges in
% black, receipts in red, the sums in pencil, a rule that holds from leaf 7 to
% the end of the volume. A transcription that records the words and the
% figures and nothing about their colour throws that field away, and the
% accounts then have to guess a receipt from the word « rec'd » --- which the
% volume stops writing on leaf 157. So the ink is recorded, per cell, as what
% is on the paper and not as what it means: \ink{red}{rec'd by check} says
% the words are red, and it is `scripts/accounts.mjs` that says red is money
% received. Black is the default and is not marked. The inks are the three
% the volume uses besides black; a fourth would be added here and in
% `scripts/render.mjs` and `scripts/tei.mjs` in the same commit.
% The file is \input, not \usepackage'd, so @ is not a letter here without saying so.
\makeatletter
\@namedef{hop@ink@red}{red!70!black}
\@namedef{hop@ink@pencil}{gray!55!black}
\@namedef{hop@ink@blue}{blue!55!black}
\newcommand{\ink}[2]{%
  \ifcsname hop@ink@#1\endcsname
    \textcolor{\csname hop@ink@#1\endcsname}{#2}%
  \else
    \PackageError{hopper}{\string\ink{#1}: unknown ink}{The inks are red, pencil and blue.}%
  \fi}
\makeatother

% A rule drawn across the money column above this row's figure.
%
% The writer rules off a run of items and writes their sum under the rule,
% and the rule is the whole of what says « this figure is a sum ». It stands
% at the head of the row the sum is on: \ruledoff & & 4355 & 00. On paper it
% is a line under the last two columns of the four a Book IV leaf has; if
% another volume ever rules a different column, the macro grows an argument.
\newcommand{\ruledoff}{\cline{3-4}}

% A diary entry begins. #1 is the date as she wrote it — « Mar. 12 »,
% « Sunday » — and #2 the date the transcriber assigns it, as a year, a month
% and a day (1947-03-12), or as much of one as the leaf allows. The first is
% what is on the paper; the second is the one editorial claim a diary needs
% that a ledger does not, because a ledger's leaf is its own address and an
% entry's date is its subject. Both are kept, apart, so the claim can be
% argued with. Used only in the notebooks; a ledger has no entries.
\newcommand{\entry}[2]{\par\medskip\noindent
  {\bfseries #1}\ifx\relax#2\relax\else\hspace{0.6em}{\footnotesize\color{gray!60!black}\textsf{#2}}\fi
  \par\nopagebreak\smallskip}

% The transcriber's note. Ours, not theirs.
\newcommand{\note}[1]{\par\smallskip{\small\color{gray!60!black}\textit{[#1]}}\par\smallskip}

% A marginal note **of theirs**, distinct from the body — and distinct from
% \note{} for exactly the reason \note{} exists.
\newcommand{\marginal}[1]{\par\smallskip\noindent\hspace{1em}%
  {\small\color{gray!60!black}#1}\par\smallskip}

% --- What is particular to this archive -------------------------------------

% Whose hand wrote it.
%
% This macro has no counterpart in the Grothendieck preamble, and it is the
% single most important thing in this one. These books were written by two
% people and annotated by more. Edward Hopper drew the work and wrote its
% title and size; Josephine Nivison Hopper wrote everything else — the
% descriptions, the anecdotes about people he refused to discuss, the prices,
% the buyers, the columns; and later hands, curatorial and unidentified, added
% to the leaves after 1967.
%
% A transcription that flattens the three has destroyed the archive's chief
% documentary value, which is that it records *two* working lives and the
% division of labour between them. #1 is `edward`, `jo`, `later` or
% `unidentified`; `unidentified` is the right answer far more often than a
% transcriber's confidence suggests, and it is not a failure to use it.
\newcommand{\hand}[2]{%
  \ifx\relax#1\relax#2\else
    \textcolor{gray!55!black}{\footnotesize\textsf{[#1]}}\,#2%
  \fi}

% Edward Hopper's ink record sketch stands here.
%
% The argument is **what is inscribed on or beside the sketch** — a title, a
% dimension, a note — and never a description of the image. Describing the
% drawing would be writing a new document: the sketch is on the site, one pane
% away, at the resolution the Whitney publishes, and prose about it competes
% with looking at it. Leave the argument empty where nothing is written.
% The mark is drawn rather than set from a symbol font: there is no
% mathematics anywhere in this archive, so no maths package is loaded, and
% $\square$ would be an undefined control sequence — which is exactly how it
% failed the first time the PDFs were compiled.
\newcommand{\sketchmark}{\raisebox{0.15ex}{\color{gray!50!black}%
  \tikz[baseline]{\draw[line width=0.4pt] (0,0) rectangle (1.1ex,1.1ex);}}}
\newcommand{\sketch}[1]{\par\smallskip\noindent
  {\small\color{gray!50!black}\sketchmark~\textsf{ink sketch}%
   \ifx\relax#1\relax\else\ --- #1\fi}\par\smallskip}

% Something printed and pasted to the leaf — a reproduction cut from a
% catalogue, a review, a dealer's label. The argument transcribes its printed
% text. These are the only places in the archive where the words are nobody's
% handwriting, and they date the leaf, which handwriting does not.
\newcommand{\clipping}[1]{\par\smallskip\noindent
  {\small\color{gray!50!black}\textsf{[clipping]}\ #1}\par\smallskip}

% A work's block opens here, under the title **exactly as the leaf gives it** —
% « Les Deux Pigeon » and « Les Poillus » stay misspelled, because the
% misspelling is Hopper's and is how the entry is found.
\newcommand{\work}[1]{\par\medskip\noindent{\bfseries #1}\par\nopagebreak\smallskip}

% --- The ruled table --------------------------------------------------------
%
% Jo Hopper ruled these columns herself and kept them for forty years; the
% table is not a presentation of the content, it *is* the content. So it is
% transcribed as a table and rendered as one, on screen and on paper, from one
% source.
%
% #1 is the column specification; #2 is the header row, `&`-separated like any
% other. A cell she left blank is left blank; a cell that cannot be read takes
% \ill{}.
%
% **Use `Y{...}` for any column that can hold a sentence** — an exhibition, a
% dealer's terms, a note — and `l`, `r` or `c` only for dates, figures and
% short names. See the note on \newcolumntype{Y} above for why that is not a
% matter of taste.
%
% The header row is emitted **bare**, between rules, and is deliberately not
% wrapped in \textbf{}. It cannot be: an `&` inside a braced group is not an
% alignment tab, so `\textbf{Date & Exhibitions}` fails at the first header
% with « Forbidden control sequence found while scanning use of
% \check@nocorr@ » — which is TeX saying, at some distance, that the tabs are
% in the wrong place. Rules carry the header instead, which is how a scholarly
% table sets one anyway. (The reading view has no such constraint and renders
% the same row as <th>.)
\newenvironment{ledgertable}[2]
  {\par\smallskip\small
   \begin{longtable}{#1}
   \hline
   #2 \\
   \hline
   \endhead}
  {\end{longtable}\par\smallskip}

% --- The appendix of notes --------------------------------------------------
%
% Everything above the appendix is on the paper. Nothing in it is.
%
% A note says what a named source says about a work these leaves record --- what
% the holding museum measures the canvas at, which other leaf carries the other
% half of a transaction --- and it is the only prose in the file that a reader
% cannot check against the photograph. On the site that is handled by putting
% the note beside a link to its source; on paper a link is not enough, because
% a PDF outlives the site it was downloaded from and travels away from it. So
% every claim here prints its source in full, name and address both, and every
% note prints the day it was written and what wrote it.
%
% The appendix is assembled by `scripts/pdf.mjs` from
% `src/content/work-notes.json` at compile time and is in no `.tex` file. That
% is the point of it. The transcription is the edition; a note is a later and
% weaker kind of writing *about* the edition, and keeping the two in one source
% file would eventually put a sentence somebody inferred inside the document
% whose whole claim on a reader is that nothing in it was inferred.
%
% Only the notes on works this batch's leaves name are printed, so the appendix
% follows the batch rather than repeating the archive, and a batch that has
% earned no notes yet gets no appendix at all --- which is a true statement
% about how much has been checked, and the same statement `work-notes.json`
% makes by being nearly empty.

% A URL is the one string in this archive that can run a hundred characters
% without a space in it, and an overfull box fails the build here because in a
% transcription an overfull box hides content. So the appendix tells TeX where
% a URL may break and sets its apparatus ragged right, rather than leaving the
% line to overrun in a file nobody recompiles.
\def\UrlBreaks{\do\/\do\-\do\.\do\_\do\?\do\&\do\=\do\+}
\Urlmuskip=0mu plus 1mu

\newcommand{\notesappendix}{\clearpage
  \par\noindent{\large\bfseries Notes on the works}\par\smallskip
  \noindent{\footnotesize\color{gray!60!black}\raggedright
    Not on the paper, and not part of the transcription. Each note below
    concerns a work named on a leaf of this batch and says what a named source
    says about it --- a museum's own catalogue record, or another leaf of these
    ledgers. Every sentence carries the source that states it, printed with its
    address so that it can be followed from a copy of this file alone. Where a
    ledger and a museum disagree the disagreement is printed rather than
    resolved.\par}\par\medskip}

% One work: the title as the leaves give it, then the note.
\newcommand{\worknote}[2]{\par\medskip\noindent{\bfseries #1}\par\nopagebreak
  \smallskip\noindent #2\par}

% One claim under it: what is asserted, and who asserts it.
\newcommand{\workclaim}[2]{\par\smallskip\noindent
  {\footnotesize\color{gray!60!black}\raggedright\hangindent=1.6em\hangafter=0
   \hspace{1em}#1\ --- #2\par}}

% Who wrote the note and when, closing the block. A note is a model's reading
% of retrieved records and is weaker evidence than anything above it; the file
% records that rather than letting the appendix pass for the edition.
\newcommand{\worknoteby}[2]{\par\smallskip\noindent
  {\scriptsize\color{gray!55!black}\hspace{1em}Written #1 by #2.\par}\par\smallskip}

% --- Title ------------------------------------------------------------------

% The watermark is drawn on the shipout background so it sits under the text of
% every page, diagonal and light enough to leave the record readable.
\AddToHook{shipout/background}{%
  \ifx\hopWatermark\empty\else
    \begin{tikzpicture}[remember picture, overlay]
      \node[rotate=52, scale=3.4, align=center, text=gray!18,
            font=\sffamily\bfseries]
        at (current page.center) {\hopWatermark};
    \end{tikzpicture}%
  \fi}

\AtBeginDocument{%
  \begin{center}
    {\large\bfseries Edward and Josephine Hopper --- \hopLedgerTitle}\\[2pt]
    \ifx\hopObject\empty\else
      {\small Whitney Museum of American Art, \hopObject}\\[4pt]
    \fi
    {\small sheets \hopFirst--\hopLast
      \ifx\hopBatch\empty\else\ (batch \hopBatch)\fi}\\[2pt]
    \ifx\hopDating\empty\else
      {\small\color{gray!60!black}The Whitney's dating: \hopDating}\\[2pt]
    \fi
    \ifx\hopWatermark\empty\else
      {\small\bfseries\color{red!45!gray}\hopWatermark}\\[2pt]
    \fi
    {\footnotesize\color{gray!60!black}Transcribed from the digitised sheets published by the
     Whitney Museum of American Art.\\
     \textcopyright{} Heirs of Josephine N. Hopper, licensed by Artists Rights Society (ARS),
     New York.\\
     Uncertain readings are underlined, illegible passages marked [\,\ldots\,],
     editorial additions bracketed.}
  \end{center}
  \vspace{1em}}

\endinput
 rather than \usepackage, and \ProvidesPackage in
% an \input-ed file makes LaTeX warn on every compile that it "requested
% package `'" and got `hopper'. A warning nobody can act on is a warning
% everybody learns to scroll past, which is how the real ones get missed.
\ProvidesFile{hopper.sty}[2026/09/05 Transcriptions of the Hopper artist's ledgers]

\usepackage{array}
% longtable, not tabularx: a leaf of thirty-one exhibition lines must break
% across pages, and tabularx cannot be wrapped in an environment of our own —
% it scans the source for a literal \end{tabularx} and dies on \end{ledgertable}
% with « File ended while scanning use of \TX@get@body ». longtable has no such
% scan, and the wrapping column type below does the job tabularx's X would.
\usepackage{longtable}

% A wrapping column, `Y{0.3}` being three tenths of the measure.
%
% Prose columns must be declared this way and not as `l`. A table of `l`
% columns is set to its natural width, which can be wider than the page, and
% TeX issues **no warning at all** because nothing ever asked the row to fit:
% the first compile of Book I produced a page whose right-hand column ran off
% the paper with a clean log. A fixed-width column wraps, and if it still
% cannot fit, TeX says so — which is what `npm run pdf` then refuses to build
% over.
%
% Leave about 0.03 of the measure per column for the inter-column space: five
% columns can share 0.85, six can share 0.80. Getting it wrong is not a
% guessing game — `npm run pdf` reports the overfull box and how many points
% too wide it is.
\newcolumntype{Y}[1]{>{\raggedright\arraybackslash}p{#1\linewidth}}
\usepackage[english]{babel}
\usepackage{fontspec}
\usepackage{geometry}
\geometry{a4paper,margin=25mm,marginparwidth=28mm}
\usepackage{marginnote}
\usepackage{xcolor}
\usepackage{tikz}
% ulem, not soul. `soul`'s \st and \ul are built on a character-by-character
% scan that XeLaTeX's font machinery breaks: the document fails with an
% undefined control sequence rather than degrading. `normalem` stops it
% hijacking \emph, which the transcriptions use for underlining of Jo Hopper's
% own.
\usepackage[normalem]{ulem}
% ulem sets each underlined word in a box TeX can neither hyphenate nor break
% after, so a paragraph dense in \uncertain{} readings can reach a state where
% no legal breakpoint exists. \emergencystretch lets TeX loosen it on a third
% pass instead of overrunning: an overfull line in a transcription hides
% content.
\setlength{\emergencystretch}{3em}
\usepackage{hyperref}
\hypersetup{colorlinks=true,linkcolor=black,urlcolor=black,pdfborder={0 0 0}}

% --- Batch metadata ---------------------------------------------------------
% Taken as they stand by the rendering script, which shows them at the head of
% the reading view. They fix what the file claims to cover: without them a
% transcription floats unmoored in an archive of five hundred sheets.

\newcommand{\ledger}[1]{\gdef\hopLedger{#1}}
\newcommand{\ledgertitle}[1]{\gdef\hopLedgerTitle{#1}}
% A notebook of Josephine Hopper's, in place of a ledger: one file to a
% notebook, filed under transcripts/notebooks/, the argument the site's slug
% (« battle-of-wash-sq »). Sets the ledger to « notebooks » for everything
% that reads it, so the title block and the watermark behave as they do for a
% batch; the rendering checks the sheets against the notebook's own listing.
\newcommand{\notebook}[1]{\gdef\hopLedger{notebooks}\gdef\hopNotebook{#1}}
\gdef\hopNotebook{}
% The Whitney's accession number for the physical volume — 96.208 … 96.213.
% This is what a citation of the object cites.
\newcommand{\objectnumber}[1]{\gdef\hopObject{#1}}
% The batch this file covers. Leaving it out drops the « (batch N) » from the
% title block rather than printing a number that would be false.
\newcommand{\batch}[1]{\gdef\hopBatch{#1}}
\newcommand{\sheets}[2]{\gdef\hopFirst{#1}\gdef\hopLast{#2}}

% The Whitney's date range for the volume, copied verbatim from its resource
% record — never inferred here, never narrowed to the years this batch
% happens to cover.
\newcommand{\dating}[1]{\gdef\hopDating{#1}}

% The legal notice, printed in the title block and repeated as a diagonal
% watermark on every page of the PDF.
%
% This is not boilerplate and it is not optional. The ledgers are © Heirs of
% Josephine N. Hopper, licensed by the Artists Rights Society; the Whitney
% records the object rights as transferred to the Museum. A transcription
% reproduces Josephine Hopper's text. Every file must therefore say what it is
% on its own face, not only in a README that does not travel with the PDF.
% A file without this line gets no watermark, so omitting it is visible in the
% output rather than silent.
\newcommand{\watermark}[1]{\gdef\hopWatermark{#1}}

\gdef\hopLedger{?}\gdef\hopLedgerTitle{}\gdef\hopObject{}\gdef\hopBatch{}%
\gdef\hopFirst{?}\gdef\hopLast{?}\gdef\hopDating{}\gdef\hopWatermark{}

% --- Keywords ---------------------------------------------------------------
%
% Closing a transcription: the vocabulary under which to search for what these
% leaves record — dealers, exhibiting societies, media, places.
% `npm run manifest` extracts them to tag the whole ledger. This line is the
% single source of the tags; there is no tags file, so no tag can describe
% leaves nobody has read.
%
% It is the one thing in the file that is not on the paper, and it earns its
% place: a reader looking for Keppel needs some way in, and a tag written by
% whoever read the sheets is the only kind that cannot describe unread ones.
\newcommand{\keywords}[1]{\par\smallskip\noindent
  {\footnotesize\textsc{Keywords} --- \textit{#1}}\par}

% --- The anchor -------------------------------------------------------------

% The sheet begins here. Two arguments, and both are needed.
%
% #1 is the Whitney's **resource ref** — 16853 — which is the sheet's whole
% address: it is what the image URL is built from and what turns the facsimile
% as the transcript is scrolled. #2 is the number **written on the paper**, or
% empty where nobody wrote one, which is the case for every cover, flyleaf,
% index and loose insertion — about one sheet in nine.
%
% Keeping both is the whole of the numbering problem in this archive. Three
% sequences overlap: the Hoppers' own leaf numbers, the Whitney's order of
% photographs, and the resource refs, which are in no order at all. Only the
% first is citable and only the third is addressable, so a transcription that
% recorded one of them would be either uncitable or unopenable.
%
% `scripts/render.mjs` checks #2 against what `scripts/catalogue.mjs` parsed
% out of the Whitney's descriptor and fails on a disagreement, because a
% mistyped leaf number silently moves an entry to another year.
\newcommand{\sheet}[2]{%
  \marginnote{\footnotesize\color{gray!70}\textsf{\ifx\relax#2\relax---\else#2\fi}}%
  \label{sheet:#1}\ignorespaces}

% --- The critical apparatus -------------------------------------------------

% Illegible: never guessed. On a ledger leaf this is most often a figure, and
% a guessed figure is a fabricated transaction.
\newcommand{\ill}{\textcolor{red!60!black}{[\,\ldots\,]}}

% A doubtful reading: the word or figure is given, and flagged as uncertain.
%
% Underlining belongs to this macro and to nothing else. Jo Hopper underlines
% her own headings --- « Etchings », « Etching Plates », « Early oil » --- and
% those are set with \emph{} instead, because a reader must be able to tell a
% doubtful reading from her emphasis at a glance and the two cannot both be
% underlines. Which of them keeps the underline is decided by consequence: an
% emphasis shown as italic loses nothing, and a doubtful figure shown as a
% certain one loses everything.
\newcommand{\uncertain}[1]{\uline{#1}}

% An editorial addition — an expanded abbreviation, an implied dollar sign.
\newcommand{\add}[1]{\textcolor{blue!50!black}{[#1]}}

% Struck out in the book. Jo Hopper struck a great deal: a price revised, a
% buyer who withdrew, an exhibition that fell through. What she crossed out is
% part of the record and often the more interesting half of it.
\newcommand{\struck}[1]{\sout{#1}}

% The ink it is written in, where the ink says something.
%
% Book IV is a state machine and the colour is its status field: charges in
% black, receipts in red, the sums in pencil, a rule that holds from leaf 7 to
% the end of the volume. A transcription that records the words and the
% figures and nothing about their colour throws that field away, and the
% accounts then have to guess a receipt from the word « rec'd » --- which the
% volume stops writing on leaf 157. So the ink is recorded, per cell, as what
% is on the paper and not as what it means: \ink{red}{rec'd by check} says
% the words are red, and it is `scripts/accounts.mjs` that says red is money
% received. Black is the default and is not marked. The inks are the three
% the volume uses besides black; a fourth would be added here and in
% `scripts/render.mjs` and `scripts/tei.mjs` in the same commit.
% The file is \input, not \usepackage'd, so @ is not a letter here without saying so.
\makeatletter
\@namedef{hop@ink@red}{red!70!black}
\@namedef{hop@ink@pencil}{gray!55!black}
\@namedef{hop@ink@blue}{blue!55!black}
\newcommand{\ink}[2]{%
  \ifcsname hop@ink@#1\endcsname
    \textcolor{\csname hop@ink@#1\endcsname}{#2}%
  \else
    \PackageError{hopper}{\string\ink{#1}: unknown ink}{The inks are red, pencil and blue.}%
  \fi}
\makeatother

% A rule drawn across the money column above this row's figure.
%
% The writer rules off a run of items and writes their sum under the rule,
% and the rule is the whole of what says « this figure is a sum ». It stands
% at the head of the row the sum is on: \ruledoff & & 4355 & 00. On paper it
% is a line under the last two columns of the four a Book IV leaf has; if
% another volume ever rules a different column, the macro grows an argument.
\newcommand{\ruledoff}{\cline{3-4}}

% A diary entry begins. #1 is the date as she wrote it — « Mar. 12 »,
% « Sunday » — and #2 the date the transcriber assigns it, as a year, a month
% and a day (1947-03-12), or as much of one as the leaf allows. The first is
% what is on the paper; the second is the one editorial claim a diary needs
% that a ledger does not, because a ledger's leaf is its own address and an
% entry's date is its subject. Both are kept, apart, so the claim can be
% argued with. Used only in the notebooks; a ledger has no entries.
\newcommand{\entry}[2]{\par\medskip\noindent
  {\bfseries #1}\ifx\relax#2\relax\else\hspace{0.6em}{\footnotesize\color{gray!60!black}\textsf{#2}}\fi
  \par\nopagebreak\smallskip}

% The transcriber's note. Ours, not theirs.
\newcommand{\note}[1]{\par\smallskip{\small\color{gray!60!black}\textit{[#1]}}\par\smallskip}

% A marginal note **of theirs**, distinct from the body — and distinct from
% \note{} for exactly the reason \note{} exists.
\newcommand{\marginal}[1]{\par\smallskip\noindent\hspace{1em}%
  {\small\color{gray!60!black}#1}\par\smallskip}

% --- What is particular to this archive -------------------------------------

% Whose hand wrote it.
%
% This macro has no counterpart in the Grothendieck preamble, and it is the
% single most important thing in this one. These books were written by two
% people and annotated by more. Edward Hopper drew the work and wrote its
% title and size; Josephine Nivison Hopper wrote everything else — the
% descriptions, the anecdotes about people he refused to discuss, the prices,
% the buyers, the columns; and later hands, curatorial and unidentified, added
% to the leaves after 1967.
%
% A transcription that flattens the three has destroyed the archive's chief
% documentary value, which is that it records *two* working lives and the
% division of labour between them. #1 is `edward`, `jo`, `later` or
% `unidentified`; `unidentified` is the right answer far more often than a
% transcriber's confidence suggests, and it is not a failure to use it.
\newcommand{\hand}[2]{%
  \ifx\relax#1\relax#2\else
    \textcolor{gray!55!black}{\footnotesize\textsf{[#1]}}\,#2%
  \fi}

% Edward Hopper's ink record sketch stands here.
%
% The argument is **what is inscribed on or beside the sketch** — a title, a
% dimension, a note — and never a description of the image. Describing the
% drawing would be writing a new document: the sketch is on the site, one pane
% away, at the resolution the Whitney publishes, and prose about it competes
% with looking at it. Leave the argument empty where nothing is written.
% The mark is drawn rather than set from a symbol font: there is no
% mathematics anywhere in this archive, so no maths package is loaded, and
% $\square$ would be an undefined control sequence — which is exactly how it
% failed the first time the PDFs were compiled.
\newcommand{\sketchmark}{\raisebox{0.15ex}{\color{gray!50!black}%
  \tikz[baseline]{\draw[line width=0.4pt] (0,0) rectangle (1.1ex,1.1ex);}}}
\newcommand{\sketch}[1]{\par\smallskip\noindent
  {\small\color{gray!50!black}\sketchmark~\textsf{ink sketch}%
   \ifx\relax#1\relax\else\ --- #1\fi}\par\smallskip}

% Something printed and pasted to the leaf — a reproduction cut from a
% catalogue, a review, a dealer's label. The argument transcribes its printed
% text. These are the only places in the archive where the words are nobody's
% handwriting, and they date the leaf, which handwriting does not.
\newcommand{\clipping}[1]{\par\smallskip\noindent
  {\small\color{gray!50!black}\textsf{[clipping]}\ #1}\par\smallskip}

% A work's block opens here, under the title **exactly as the leaf gives it** —
% « Les Deux Pigeon » and « Les Poillus » stay misspelled, because the
% misspelling is Hopper's and is how the entry is found.
\newcommand{\work}[1]{\par\medskip\noindent{\bfseries #1}\par\nopagebreak\smallskip}

% --- The ruled table --------------------------------------------------------
%
% Jo Hopper ruled these columns herself and kept them for forty years; the
% table is not a presentation of the content, it *is* the content. So it is
% transcribed as a table and rendered as one, on screen and on paper, from one
% source.
%
% #1 is the column specification; #2 is the header row, `&`-separated like any
% other. A cell she left blank is left blank; a cell that cannot be read takes
% \ill{}.
%
% **Use `Y{...}` for any column that can hold a sentence** — an exhibition, a
% dealer's terms, a note — and `l`, `r` or `c` only for dates, figures and
% short names. See the note on \newcolumntype{Y} above for why that is not a
% matter of taste.
%
% The header row is emitted **bare**, between rules, and is deliberately not
% wrapped in \textbf{}. It cannot be: an `&` inside a braced group is not an
% alignment tab, so `\textbf{Date & Exhibitions}` fails at the first header
% with « Forbidden control sequence found while scanning use of
% \check@nocorr@ » — which is TeX saying, at some distance, that the tabs are
% in the wrong place. Rules carry the header instead, which is how a scholarly
% table sets one anyway. (The reading view has no such constraint and renders
% the same row as <th>.)
\newenvironment{ledgertable}[2]
  {\par\smallskip\small
   \begin{longtable}{#1}
   \hline
   #2 \\
   \hline
   \endhead}
  {\end{longtable}\par\smallskip}

% --- The appendix of notes --------------------------------------------------
%
% Everything above the appendix is on the paper. Nothing in it is.
%
% A note says what a named source says about a work these leaves record --- what
% the holding museum measures the canvas at, which other leaf carries the other
% half of a transaction --- and it is the only prose in the file that a reader
% cannot check against the photograph. On the site that is handled by putting
% the note beside a link to its source; on paper a link is not enough, because
% a PDF outlives the site it was downloaded from and travels away from it. So
% every claim here prints its source in full, name and address both, and every
% note prints the day it was written and what wrote it.
%
% The appendix is assembled by `scripts/pdf.mjs` from
% `src/content/work-notes.json` at compile time and is in no `.tex` file. That
% is the point of it. The transcription is the edition; a note is a later and
% weaker kind of writing *about* the edition, and keeping the two in one source
% file would eventually put a sentence somebody inferred inside the document
% whose whole claim on a reader is that nothing in it was inferred.
%
% Only the notes on works this batch's leaves name are printed, so the appendix
% follows the batch rather than repeating the archive, and a batch that has
% earned no notes yet gets no appendix at all --- which is a true statement
% about how much has been checked, and the same statement `work-notes.json`
% makes by being nearly empty.

% A URL is the one string in this archive that can run a hundred characters
% without a space in it, and an overfull box fails the build here because in a
% transcription an overfull box hides content. So the appendix tells TeX where
% a URL may break and sets its apparatus ragged right, rather than leaving the
% line to overrun in a file nobody recompiles.
\def\UrlBreaks{\do\/\do\-\do\.\do\_\do\?\do\&\do\=\do\+}
\Urlmuskip=0mu plus 1mu

\newcommand{\notesappendix}{\clearpage
  \par\noindent{\large\bfseries Notes on the works}\par\smallskip
  \noindent{\footnotesize\color{gray!60!black}\raggedright
    Not on the paper, and not part of the transcription. Each note below
    concerns a work named on a leaf of this batch and says what a named source
    says about it --- a museum's own catalogue record, or another leaf of these
    ledgers. Every sentence carries the source that states it, printed with its
    address so that it can be followed from a copy of this file alone. Where a
    ledger and a museum disagree the disagreement is printed rather than
    resolved.\par}\par\medskip}

% One work: the title as the leaves give it, then the note.
\newcommand{\worknote}[2]{\par\medskip\noindent{\bfseries #1}\par\nopagebreak
  \smallskip\noindent #2\par}

% One claim under it: what is asserted, and who asserts it.
\newcommand{\workclaim}[2]{\par\smallskip\noindent
  {\footnotesize\color{gray!60!black}\raggedright\hangindent=1.6em\hangafter=0
   \hspace{1em}#1\ --- #2\par}}

% Who wrote the note and when, closing the block. A note is a model's reading
% of retrieved records and is weaker evidence than anything above it; the file
% records that rather than letting the appendix pass for the edition.
\newcommand{\worknoteby}[2]{\par\smallskip\noindent
  {\scriptsize\color{gray!55!black}\hspace{1em}Written #1 by #2.\par}\par\smallskip}

% --- Title ------------------------------------------------------------------

% The watermark is drawn on the shipout background so it sits under the text of
% every page, diagonal and light enough to leave the record readable.
\AddToHook{shipout/background}{%
  \ifx\hopWatermark\empty\else
    \begin{tikzpicture}[remember picture, overlay]
      \node[rotate=52, scale=3.4, align=center, text=gray!18,
            font=\sffamily\bfseries]
        at (current page.center) {\hopWatermark};
    \end{tikzpicture}%
  \fi}

\AtBeginDocument{%
  \begin{center}
    {\large\bfseries Edward and Josephine Hopper --- \hopLedgerTitle}\\[2pt]
    \ifx\hopObject\empty\else
      {\small Whitney Museum of American Art, \hopObject}\\[4pt]
    \fi
    {\small sheets \hopFirst--\hopLast
      \ifx\hopBatch\empty\else\ (batch \hopBatch)\fi}\\[2pt]
    \ifx\hopDating\empty\else
      {\small\color{gray!60!black}The Whitney's dating: \hopDating}\\[2pt]
    \fi
    \ifx\hopWatermark\empty\else
      {\small\bfseries\color{red!45!gray}\hopWatermark}\\[2pt]
    \fi
    {\footnotesize\color{gray!60!black}Transcribed from the digitised sheets published by the
     Whitney Museum of American Art.\\
     \textcopyright{} Heirs of Josephine N. Hopper, licensed by Artists Rights Society (ARS),
     New York.\\
     Uncertain readings are underlined, illegible passages marked [\,\ldots\,],
     editorial additions bracketed.}
  \end{center}
  \vspace{1em}}

\endinput
 rather than \usepackage, and \ProvidesPackage in
% an \input-ed file makes LaTeX warn on every compile that it "requested
% package `'" and got `hopper'. A warning nobody can act on is a warning
% everybody learns to scroll past, which is how the real ones get missed.
\ProvidesFile{hopper.sty}[2026/09/05 Transcriptions of the Hopper artist's ledgers]

\usepackage{array}
% longtable, not tabularx: a leaf of thirty-one exhibition lines must break
% across pages, and tabularx cannot be wrapped in an environment of our own —
% it scans the source for a literal \end{tabularx} and dies on \end{ledgertable}
% with « File ended while scanning use of \TX@get@body ». longtable has no such
% scan, and the wrapping column type below does the job tabularx's X would.
\usepackage{longtable}

% A wrapping column, `Y{0.3}` being three tenths of the measure.
%
% Prose columns must be declared this way and not as `l`. A table of `l`
% columns is set to its natural width, which can be wider than the page, and
% TeX issues **no warning at all** because nothing ever asked the row to fit:
% the first compile of Book I produced a page whose right-hand column ran off
% the paper with a clean log. A fixed-width column wraps, and if it still
% cannot fit, TeX says so — which is what `npm run pdf` then refuses to build
% over.
%
% Leave about 0.03 of the measure per column for the inter-column space: five
% columns can share 0.85, six can share 0.80. Getting it wrong is not a
% guessing game — `npm run pdf` reports the overfull box and how many points
% too wide it is.
\newcolumntype{Y}[1]{>{\raggedright\arraybackslash}p{#1\linewidth}}
\usepackage[english]{babel}
\usepackage{fontspec}
\usepackage{geometry}
\geometry{a4paper,margin=25mm,marginparwidth=28mm}
\usepackage{marginnote}
\usepackage{xcolor}
\usepackage{tikz}
% ulem, not soul. `soul`'s \st and \ul are built on a character-by-character
% scan that XeLaTeX's font machinery breaks: the document fails with an
% undefined control sequence rather than degrading. `normalem` stops it
% hijacking \emph, which the transcriptions use for underlining of Jo Hopper's
% own.
\usepackage[normalem]{ulem}
% ulem sets each underlined word in a box TeX can neither hyphenate nor break
% after, so a paragraph dense in \uncertain{} readings can reach a state where
% no legal breakpoint exists. \emergencystretch lets TeX loosen it on a third
% pass instead of overrunning: an overfull line in a transcription hides
% content.
\setlength{\emergencystretch}{3em}
\usepackage{hyperref}
\hypersetup{colorlinks=true,linkcolor=black,urlcolor=black,pdfborder={0 0 0}}

% --- Batch metadata ---------------------------------------------------------
% Taken as they stand by the rendering script, which shows them at the head of
% the reading view. They fix what the file claims to cover: without them a
% transcription floats unmoored in an archive of five hundred sheets.

\newcommand{\ledger}[1]{\gdef\hopLedger{#1}}
\newcommand{\ledgertitle}[1]{\gdef\hopLedgerTitle{#1}}
% A notebook of Josephine Hopper's, in place of a ledger: one file to a
% notebook, filed under transcripts/notebooks/, the argument the site's slug
% (« battle-of-wash-sq »). Sets the ledger to « notebooks » for everything
% that reads it, so the title block and the watermark behave as they do for a
% batch; the rendering checks the sheets against the notebook's own listing.
\newcommand{\notebook}[1]{\gdef\hopLedger{notebooks}\gdef\hopNotebook{#1}}
\gdef\hopNotebook{}
% The Whitney's accession number for the physical volume — 96.208 … 96.213.
% This is what a citation of the object cites.
\newcommand{\objectnumber}[1]{\gdef\hopObject{#1}}
% The batch this file covers. Leaving it out drops the « (batch N) » from the
% title block rather than printing a number that would be false.
\newcommand{\batch}[1]{\gdef\hopBatch{#1}}
\newcommand{\sheets}[2]{\gdef\hopFirst{#1}\gdef\hopLast{#2}}

% The Whitney's date range for the volume, copied verbatim from its resource
% record — never inferred here, never narrowed to the years this batch
% happens to cover.
\newcommand{\dating}[1]{\gdef\hopDating{#1}}

% The legal notice, printed in the title block and repeated as a diagonal
% watermark on every page of the PDF.
%
% This is not boilerplate and it is not optional. The ledgers are © Heirs of
% Josephine N. Hopper, licensed by the Artists Rights Society; the Whitney
% records the object rights as transferred to the Museum. A transcription
% reproduces Josephine Hopper's text. Every file must therefore say what it is
% on its own face, not only in a README that does not travel with the PDF.
% A file without this line gets no watermark, so omitting it is visible in the
% output rather than silent.
\newcommand{\watermark}[1]{\gdef\hopWatermark{#1}}

\gdef\hopLedger{?}\gdef\hopLedgerTitle{}\gdef\hopObject{}\gdef\hopBatch{}%
\gdef\hopFirst{?}\gdef\hopLast{?}\gdef\hopDating{}\gdef\hopWatermark{}

% --- Keywords ---------------------------------------------------------------
%
% Closing a transcription: the vocabulary under which to search for what these
% leaves record — dealers, exhibiting societies, media, places.
% `npm run manifest` extracts them to tag the whole ledger. This line is the
% single source of the tags; there is no tags file, so no tag can describe
% leaves nobody has read.
%
% It is the one thing in the file that is not on the paper, and it earns its
% place: a reader looking for Keppel needs some way in, and a tag written by
% whoever read the sheets is the only kind that cannot describe unread ones.
\newcommand{\keywords}[1]{\par\smallskip\noindent
  {\footnotesize\textsc{Keywords} --- \textit{#1}}\par}

% --- The anchor -------------------------------------------------------------

% The sheet begins here. Two arguments, and both are needed.
%
% #1 is the Whitney's **resource ref** — 16853 — which is the sheet's whole
% address: it is what the image URL is built from and what turns the facsimile
% as the transcript is scrolled. #2 is the number **written on the paper**, or
% empty where nobody wrote one, which is the case for every cover, flyleaf,
% index and loose insertion — about one sheet in nine.
%
% Keeping both is the whole of the numbering problem in this archive. Three
% sequences overlap: the Hoppers' own leaf numbers, the Whitney's order of
% photographs, and the resource refs, which are in no order at all. Only the
% first is citable and only the third is addressable, so a transcription that
% recorded one of them would be either uncitable or unopenable.
%
% `scripts/render.mjs` checks #2 against what `scripts/catalogue.mjs` parsed
% out of the Whitney's descriptor and fails on a disagreement, because a
% mistyped leaf number silently moves an entry to another year.
\newcommand{\sheet}[2]{%
  \marginnote{\footnotesize\color{gray!70}\textsf{\ifx\relax#2\relax---\else#2\fi}}%
  \label{sheet:#1}\ignorespaces}

% --- The critical apparatus -------------------------------------------------

% Illegible: never guessed. On a ledger leaf this is most often a figure, and
% a guessed figure is a fabricated transaction.
\newcommand{\ill}{\textcolor{red!60!black}{[\,\ldots\,]}}

% A doubtful reading: the word or figure is given, and flagged as uncertain.
%
% Underlining belongs to this macro and to nothing else. Jo Hopper underlines
% her own headings --- « Etchings », « Etching Plates », « Early oil » --- and
% those are set with \emph{} instead, because a reader must be able to tell a
% doubtful reading from her emphasis at a glance and the two cannot both be
% underlines. Which of them keeps the underline is decided by consequence: an
% emphasis shown as italic loses nothing, and a doubtful figure shown as a
% certain one loses everything.
\newcommand{\uncertain}[1]{\uline{#1}}

% An editorial addition — an expanded abbreviation, an implied dollar sign.
\newcommand{\add}[1]{\textcolor{blue!50!black}{[#1]}}

% Struck out in the book. Jo Hopper struck a great deal: a price revised, a
% buyer who withdrew, an exhibition that fell through. What she crossed out is
% part of the record and often the more interesting half of it.
\newcommand{\struck}[1]{\sout{#1}}

% The ink it is written in, where the ink says something.
%
% Book IV is a state machine and the colour is its status field: charges in
% black, receipts in red, the sums in pencil, a rule that holds from leaf 7 to
% the end of the volume. A transcription that records the words and the
% figures and nothing about their colour throws that field away, and the
% accounts then have to guess a receipt from the word « rec'd » --- which the
% volume stops writing on leaf 157. So the ink is recorded, per cell, as what
% is on the paper and not as what it means: \ink{red}{rec'd by check} says
% the words are red, and it is `scripts/accounts.mjs` that says red is money
% received. Black is the default and is not marked. The inks are the three
% the volume uses besides black; a fourth would be added here and in
% `scripts/render.mjs` and `scripts/tei.mjs` in the same commit.
% The file is \input, not \usepackage'd, so @ is not a letter here without saying so.
\makeatletter
\@namedef{hop@ink@red}{red!70!black}
\@namedef{hop@ink@pencil}{gray!55!black}
\@namedef{hop@ink@blue}{blue!55!black}
\newcommand{\ink}[2]{%
  \ifcsname hop@ink@#1\endcsname
    \textcolor{\csname hop@ink@#1\endcsname}{#2}%
  \else
    \PackageError{hopper}{\string\ink{#1}: unknown ink}{The inks are red, pencil and blue.}%
  \fi}
\makeatother

% A rule drawn across the money column above this row's figure.
%
% The writer rules off a run of items and writes their sum under the rule,
% and the rule is the whole of what says « this figure is a sum ». It stands
% at the head of the row the sum is on: \ruledoff & & 4355 & 00. On paper it
% is a line under the last two columns of the four a Book IV leaf has; if
% another volume ever rules a different column, the macro grows an argument.
\newcommand{\ruledoff}{\cline{3-4}}

% A diary entry begins. #1 is the date as she wrote it — « Mar. 12 »,
% « Sunday » — and #2 the date the transcriber assigns it, as a year, a month
% and a day (1947-03-12), or as much of one as the leaf allows. The first is
% what is on the paper; the second is the one editorial claim a diary needs
% that a ledger does not, because a ledger's leaf is its own address and an
% entry's date is its subject. Both are kept, apart, so the claim can be
% argued with. Used only in the notebooks; a ledger has no entries.
\newcommand{\entry}[2]{\par\medskip\noindent
  {\bfseries #1}\ifx\relax#2\relax\else\hspace{0.6em}{\footnotesize\color{gray!60!black}\textsf{#2}}\fi
  \par\nopagebreak\smallskip}

% The transcriber's note. Ours, not theirs.
\newcommand{\note}[1]{\par\smallskip{\small\color{gray!60!black}\textit{[#1]}}\par\smallskip}

% A marginal note **of theirs**, distinct from the body — and distinct from
% \note{} for exactly the reason \note{} exists.
\newcommand{\marginal}[1]{\par\smallskip\noindent\hspace{1em}%
  {\small\color{gray!60!black}#1}\par\smallskip}

% --- What is particular to this archive -------------------------------------

% Whose hand wrote it.
%
% This macro has no counterpart in the Grothendieck preamble, and it is the
% single most important thing in this one. These books were written by two
% people and annotated by more. Edward Hopper drew the work and wrote its
% title and size; Josephine Nivison Hopper wrote everything else — the
% descriptions, the anecdotes about people he refused to discuss, the prices,
% the buyers, the columns; and later hands, curatorial and unidentified, added
% to the leaves after 1967.
%
% A transcription that flattens the three has destroyed the archive's chief
% documentary value, which is that it records *two* working lives and the
% division of labour between them. #1 is `edward`, `jo`, `later` or
% `unidentified`; `unidentified` is the right answer far more often than a
% transcriber's confidence suggests, and it is not a failure to use it.
\newcommand{\hand}[2]{%
  \ifx\relax#1\relax#2\else
    \textcolor{gray!55!black}{\footnotesize\textsf{[#1]}}\,#2%
  \fi}

% Edward Hopper's ink record sketch stands here.
%
% The argument is **what is inscribed on or beside the sketch** — a title, a
% dimension, a note — and never a description of the image. Describing the
% drawing would be writing a new document: the sketch is on the site, one pane
% away, at the resolution the Whitney publishes, and prose about it competes
% with looking at it. Leave the argument empty where nothing is written.
% The mark is drawn rather than set from a symbol font: there is no
% mathematics anywhere in this archive, so no maths package is loaded, and
% $\square$ would be an undefined control sequence — which is exactly how it
% failed the first time the PDFs were compiled.
\newcommand{\sketchmark}{\raisebox{0.15ex}{\color{gray!50!black}%
  \tikz[baseline]{\draw[line width=0.4pt] (0,0) rectangle (1.1ex,1.1ex);}}}
\newcommand{\sketch}[1]{\par\smallskip\noindent
  {\small\color{gray!50!black}\sketchmark~\textsf{ink sketch}%
   \ifx\relax#1\relax\else\ --- #1\fi}\par\smallskip}

% Something printed and pasted to the leaf — a reproduction cut from a
% catalogue, a review, a dealer's label. The argument transcribes its printed
% text. These are the only places in the archive where the words are nobody's
% handwriting, and they date the leaf, which handwriting does not.
\newcommand{\clipping}[1]{\par\smallskip\noindent
  {\small\color{gray!50!black}\textsf{[clipping]}\ #1}\par\smallskip}

% A work's block opens here, under the title **exactly as the leaf gives it** —
% « Les Deux Pigeon » and « Les Poillus » stay misspelled, because the
% misspelling is Hopper's and is how the entry is found.
\newcommand{\work}[1]{\par\medskip\noindent{\bfseries #1}\par\nopagebreak\smallskip}

% --- The ruled table --------------------------------------------------------
%
% Jo Hopper ruled these columns herself and kept them for forty years; the
% table is not a presentation of the content, it *is* the content. So it is
% transcribed as a table and rendered as one, on screen and on paper, from one
% source.
%
% #1 is the column specification; #2 is the header row, `&`-separated like any
% other. A cell she left blank is left blank; a cell that cannot be read takes
% \ill{}.
%
% **Use `Y{...}` for any column that can hold a sentence** — an exhibition, a
% dealer's terms, a note — and `l`, `r` or `c` only for dates, figures and
% short names. See the note on \newcolumntype{Y} above for why that is not a
% matter of taste.
%
% The header row is emitted **bare**, between rules, and is deliberately not
% wrapped in \textbf{}. It cannot be: an `&` inside a braced group is not an
% alignment tab, so `\textbf{Date & Exhibitions}` fails at the first header
% with « Forbidden control sequence found while scanning use of
% \check@nocorr@ » — which is TeX saying, at some distance, that the tabs are
% in the wrong place. Rules carry the header instead, which is how a scholarly
% table sets one anyway. (The reading view has no such constraint and renders
% the same row as <th>.)
\newenvironment{ledgertable}[2]
  {\par\smallskip\small
   \begin{longtable}{#1}
   \hline
   #2 \\
   \hline
   \endhead}
  {\end{longtable}\par\smallskip}

% --- The appendix of notes --------------------------------------------------
%
% Everything above the appendix is on the paper. Nothing in it is.
%
% A note says what a named source says about a work these leaves record --- what
% the holding museum measures the canvas at, which other leaf carries the other
% half of a transaction --- and it is the only prose in the file that a reader
% cannot check against the photograph. On the site that is handled by putting
% the note beside a link to its source; on paper a link is not enough, because
% a PDF outlives the site it was downloaded from and travels away from it. So
% every claim here prints its source in full, name and address both, and every
% note prints the day it was written and what wrote it.
%
% The appendix is assembled by `scripts/pdf.mjs` from
% `src/content/work-notes.json` at compile time and is in no `.tex` file. That
% is the point of it. The transcription is the edition; a note is a later and
% weaker kind of writing *about* the edition, and keeping the two in one source
% file would eventually put a sentence somebody inferred inside the document
% whose whole claim on a reader is that nothing in it was inferred.
%
% Only the notes on works this batch's leaves name are printed, so the appendix
% follows the batch rather than repeating the archive, and a batch that has
% earned no notes yet gets no appendix at all --- which is a true statement
% about how much has been checked, and the same statement `work-notes.json`
% makes by being nearly empty.

% A URL is the one string in this archive that can run a hundred characters
% without a space in it, and an overfull box fails the build here because in a
% transcription an overfull box hides content. So the appendix tells TeX where
% a URL may break and sets its apparatus ragged right, rather than leaving the
% line to overrun in a file nobody recompiles.
\def\UrlBreaks{\do\/\do\-\do\.\do\_\do\?\do\&\do\=\do\+}
\Urlmuskip=0mu plus 1mu

\newcommand{\notesappendix}{\clearpage
  \par\noindent{\large\bfseries Notes on the works}\par\smallskip
  \noindent{\footnotesize\color{gray!60!black}\raggedright
    Not on the paper, and not part of the transcription. Each note below
    concerns a work named on a leaf of this batch and says what a named source
    says about it --- a museum's own catalogue record, or another leaf of these
    ledgers. Every sentence carries the source that states it, printed with its
    address so that it can be followed from a copy of this file alone. Where a
    ledger and a museum disagree the disagreement is printed rather than
    resolved.\par}\par\medskip}

% One work: the title as the leaves give it, then the note.
\newcommand{\worknote}[2]{\par\medskip\noindent{\bfseries #1}\par\nopagebreak
  \smallskip\noindent #2\par}

% One claim under it: what is asserted, and who asserts it.
\newcommand{\workclaim}[2]{\par\smallskip\noindent
  {\footnotesize\color{gray!60!black}\raggedright\hangindent=1.6em\hangafter=0
   \hspace{1em}#1\ --- #2\par}}

% Who wrote the note and when, closing the block. A note is a model's reading
% of retrieved records and is weaker evidence than anything above it; the file
% records that rather than letting the appendix pass for the edition.
\newcommand{\worknoteby}[2]{\par\smallskip\noindent
  {\scriptsize\color{gray!55!black}\hspace{1em}Written #1 by #2.\par}\par\smallskip}

% --- Title ------------------------------------------------------------------

% The watermark is drawn on the shipout background so it sits under the text of
% every page, diagonal and light enough to leave the record readable.
\AddToHook{shipout/background}{%
  \ifx\hopWatermark\empty\else
    \begin{tikzpicture}[remember picture, overlay]
      \node[rotate=52, scale=3.4, align=center, text=gray!18,
            font=\sffamily\bfseries]
        at (current page.center) {\hopWatermark};
    \end{tikzpicture}%
  \fi}

\AtBeginDocument{%
  \begin{center}
    {\large\bfseries Edward and Josephine Hopper --- \hopLedgerTitle}\\[2pt]
    \ifx\hopObject\empty\else
      {\small Whitney Museum of American Art, \hopObject}\\[4pt]
    \fi
    {\small sheets \hopFirst--\hopLast
      \ifx\hopBatch\empty\else\ (batch \hopBatch)\fi}\\[2pt]
    \ifx\hopDating\empty\else
      {\small\color{gray!60!black}The Whitney's dating: \hopDating}\\[2pt]
    \fi
    \ifx\hopWatermark\empty\else
      {\small\bfseries\color{red!45!gray}\hopWatermark}\\[2pt]
    \fi
    {\footnotesize\color{gray!60!black}Transcribed from the digitised sheets published by the
     Whitney Museum of American Art.\\
     \textcopyright{} Heirs of Josephine N. Hopper, licensed by Artists Rights Society (ARS),
     New York.\\
     Uncertain readings are underlined, illegible passages marked [\,\ldots\,],
     editorial additions bracketed.}
  \end{center}
  \vspace{1em}}

\endinput


\ledger{book-i}
\ledgertitle{Artist's ledger --- Book I}
\objectnumber{96.208}
\batch{7}
\sheets{73}{84}
\dating{1913--1963}
\watermark{Demonstration edition}

\begin{document}

\section{Leaf 64: Water Colors continued, 1926}

\sheet{16704}{64}

\hand{jo}{1926 Trip train to Eastport, Me; then to Bangor, on by boat to
Rockland - then ± 7 weeks then on to Gloucester.}

\hand{jo}{\emph{Water Colors}, Con. \qquad 1926}

\begin{ledgertable}{Y{0.10}Y{0.06}Y{0.44}Y{0.09}Y{0.14}}{Delivered to Rehn & & & Amt. rec'd. & When rec'd.}
Aug. 15, 26 & & Bow of Osprey - deck, ugly tear chaffing gear, rust color \quad 14 x 20 \quad Rockland \quad Frank Crowninshield. Ap. 1928 \quad 275 - 1/3 com. & 183.33 & June 8, 28. \\
" & 250 & Lime Rock Quarry II (Rockland) \quad 250 - 1/3 com. \quad Mrs. Blanchard \quad Oct. 27. \quad 14 x 20 & 166.66 & Nov. 4, '27. \\
" & & Davis House, Middle St., yellow, sunlight on house + tree \quad 14 x 20 \quad \uncertain{Somerville} \quad \uncertain{Harry Hessing} (1926) \quad 200 - 1/3 & 133 1/3 & Mar. 28, 1930. \\
" & 300 & Gloucester Harbor - rocks, breakwater effect? far shore sharply in detail - not very unique composition. \quad 14 x 20 \quad Utica Show - to friend of Ed. Root, Mar. 1928. Mrs. Kellog. & 200 & June 8, '28. \\
Sept. 30, 26 & & Schorner's Hall - big flat floor, crowded sheds + buildings surrounding it, grey sky, Rockland. \quad 14 x 20 \quad Rob't. Wheelwright of Phil. \quad owns Pemaquid Light p. 72. & 200 & Mar. 9, 28. \\
" & & Lime Rock Quarry I - darker than II \quad 14 x 20 \quad Rockland. \quad Chas. E. Buckley. \quad 800 - 1/3 = 533 1/3 & & Dec. 30, 1958. \\
" & 250. & trawler + telegraph pole with spikes to climb \quad (profile) deep, below docks. \quad 14 x 20 (Rockland.) \quad 250 - 1/3 com. = & 166 2/3 & Mar. 4, 27. \\
 & & bgt. by Dr. Hall of Princetown & & \\
" & 250 & Rockland Harbor - bright red house, white buildings, blue \quad big water + sky, pile of red bricks R. lower corner. \quad 14 x 20 \quad Dr. T. L. Bennett \quad May 19, 27 \quad Bennet auction, fall 1927, bgt. in by Rehn for \$120? \quad Resold to Hartford Antheneum - own Capt. Stout's House p. 66. & 166 2/3 & May 19, 27 \\
" & & Ship Yard Rockland - long side of boat (Butterfly) + roofs beyond \quad 14 x 20 \quad Feb. 1930 - Yankee living in Chili. Wheelwright owns Pemaquid Light + Rockland Harbor, above: could it be an exchange for Pemaquid Light? & & \\
" & 250 & Schorner's Bowsprit - derelict drawn ashore, glorious \quad 14 x 20 \quad Rockland \quad Mr. C. B. Hoyt (of Fog Museum counsel) Jan. 27. \quad Mar. 4, 27. & 166 2/3 & Mar. 4, 27. \\
" \quad Rockland & & Lime Rock R.R. \quad 14 x 20 \quad White house, shingles, dark tree + tracks across foreground \quad Dr. Hall. (Princeton) & 166 2/3 & Mar. 9, 28. \\
" \quad Rockland & & Beam trawler Osprey \quad 14 x 20 \quad delicate profile + white \uncertain{slim} dock foreground (Rockland) \quad Steinberg \quad 750 - 1/3 = 500 & & July 8, '57. \\
" \quad Rockland & & Beam trawler Widgeon \quad 14 x 20 \quad chocolate profile + rigging, far shore sharp. \quad Bgt. by Adolph Lewisohn - Jan. 27. \quad \ill{} & 166 2/3 & Ap. 23, 27. \\
\end{ledgertable}

\marginal{See Bk. IV \quad --- beside the Lime Rock Quarry I row}

\marginal{Returned to Studio Nov. 12, 1936 \quad --- written up the left margin
beside the Ship Yard Rockland row}

\marginal{Sent from Gloucester to Rehn Gallery.}

\marginal{Book IV \quad due for \quad --- beside the Beam trawler Widgeon row}

\note{« Rockland » in the date column of the last four rows is in blue crayon
and later than the ink. « Somerville » in the Davis House row is in the same
blue crayon and neither it nor « Harry Hessing » on the line below will settle:
the medial letters of both are the ones this hand does not distinguish. So is
the word after « white » in the Beam trawler Osprey description, offered here
as « slim ».}

\note{the \ill{} at the end of the last row is a short word written under
« Adolph Lewisohn » at the very foot of the leaf, cut by the edge of the
photograph; only the tops of its letters survive. « Fog Museum » is as written,
for the Fogg, confirming leaf 62. The size « 14 x 20 » is written above the
line rather than in it on most rows of this leaf and is given in the cell.}

\note{a small pencil « 133 1/3 » with an arrow drawn up to the row above stands
below the Kellog line. The arrow is not interpreted here. The right-hand edge
of this photograph shows the facing leaf 65 at the fore-edge; its column is
transcribed at sheet 16448.}

\section{Loose photograph between leaves 64 and 65: Manhattan Bridge Entrance}

\sheet{18051}{}

\note{a loose photograph of the Manhattan Bridge Entrance watercolour, laid
between leaves 64 and 65. Nothing is written on the face.}

\section{Loose photograph, verso}

\sheet{17242}{}

\hand{jo}{Manhattan Bridge Entrance\\
Spring of 1926.\\
Owned by Prof. A. Percy Saunders\\
Hamilton College, Clinton, N.Y.\\
Jan. 1928.}

\note{the picture is the last entry of leaf 65, which records the same buyer
and the same month.}

\section{Leaf 65: Water Colors continued, 1926--27}

\sheet{16448}{65}

\hand{jo}{\emph{Water Colors} - Con. \qquad 1926, '27.}

\marginal{Rockland 1926.}

\begin{ledgertable}{Y{0.10}Y{0.06}Y{0.44}Y{0.09}Y{0.14}}{Delivered to Rehn & Price & & Amt. & When rec'd}
Sept. 30, 26 & \$250. & Haunted House \quad 14 x 20 \quad old boardedup boarding house, Rockland \quad near ship yard. \quad Mr. Leslie G. Shaffer, early Feb. & \$166 2/3 & Mar. 4, 27. \\
" & & Civil War Camp Ground \quad 14 x 20 \quad telegraph pole, straw colored fields, Rockland - Taken in exchange \quad Taken in exchange for 2 Lights Village, Dr. Lisle. & 200 & Mar. 9, '28. \\
" & & Harbor Shore, Rockland - beach, rocks, house above. \quad crescent sweep of shore - green grass & & \\
" & & R.R. Crossing - empty fair grounds, Rockland \quad 1926. \quad Hubert A. (?) Goldstone \quad 600 - 1/3 = 400 \quad Baker photographs \quad who now owns Moonlight Interior oil. & & July 8, 1957. \\
" & & Trees. East Gloucester. Wild + woolly. 12 x 18 \quad See Book IV \quad Alfredo Valente \quad 2000 - 1/3 = & 1333 1/3 & June 25 1962. \\
" & 250 & Universalist Church tower beyond roofs + gables. \quad 14 x 20 \quad Gloucester \quad Dr. Hall \quad Bgt. by Dr. T. L. Bennett. At Bennet Auction, fall 1927, bgt. in by Rehn for 300. Resold to Dr. Hall. & 166 2/3 & May 19, 27. \\
" & \$250 & The Hill - Gloucester St. looking down, shadows crossing st. \quad 14 x 20 \quad Mr. Briggs, Boston & \$166 2/3 & Mar. 4, 27. \\
Feb. 7, 27 & & Tony's House - bright blue sky, white house, stone wall - \quad 14 x 20 \quad happy one (where roosters lived) got its spot. \quad Gloucester \quad Rehn - ? \quad Bgt. by Rehn? or which other one if not this? Paid for. \quad Mr. + Mrs. Seymour H. \quad Albright Gal., Buffalo, N.Y. \quad , Knox \quad 300 - 1/3 & 200. & Mar. 15, 1930. \\
" & 250 & Abbot's House - cold day Oct. \quad 14 x 20 \quad Gloucester \quad Bgt. by Dr. Hall - Feb. 27. Returned Jan. 1929. \quad Bgt. by New Britain Mus. from Rehn Gal. Mar. 1946. & 166 2/3 & May 18, 27 \\
Aug. 15, 26 & 250 & Anderson's house - darkish, strong. Bay window \quad 14 x 20 \quad Gloucester \quad Bgt. by John T. \struck{Spaulding}, Feb. \quad Boston Mus. 1948. & 166 2/3 & Ap. 23, 27. \\
Feb. 7, 27 & \$250 & House by Squam River, rear, light on back of house, door of shed. Rive. across front \quad Bgt. by Mr. John Spalding, Feb. \quad 14 x 20 \quad Boston Mus. 1948. & 166 2/3 & Ap. 23, 27. \\
fall 26 & & Manhattan Bridge Entrance - big arch, beauty. \quad 14 x 20. \quad Prof. A. P. Saunders, Hamilton College, Clinton, N.Y. \quad Jan. 1928. & 200. & June 8, 28. \\
\end{ledgertable}

\marginal{See Bk. IV. \quad --- beside the Harbor Shore row}

\marginal{to Rehn Gallery \quad Painted 1926. Gloucester \quad --- above and
below the date of the House by Squam River row}

\marginal{Painted in Spring \quad --- above « fall 26 »}

\marginal{Painted in N.Y. Studio}

\note{a single large brace at the left gathers the Civil War Camp Ground,
Harbor Shore and R.R. Crossing rows, which share the date above them. The
« (?) » after « Hubert A. » is hers, written small and ringed above the line,
and is not an editorial mark. « , Knox » in the Tony's House row is written on
the line below « Mr. + Mrs. Seymour H. », with « Albright Gal., Buffalo, N.Y. »
continuing the line above it; the two lines are given in the order they stand.}

\note{« Painted in N.Y. Studio » is in blue crayon at the foot of the leaf,
below the last row, and is not part of it. « Spaulding » in the Anderson's
house row is struck and rewritten « Spalding » two lines below.}

\section{Leaf 66: Watercolors 1927 (clipping down)}

\sheet{17992}{66}

\note{leaf 66 photographed with the Custom House, Portland halftone lying over
the lower half of the leaf, where it covers eight of the twelve entries. The
transcription of this leaf is at sheet 16820, which shows them all; nothing is
hidden there that this photograph shows.}

\clipping{a halftone of a watercolour --- a pillared front, windows, wet
pavement --- hinged at its upper edge and lying over the lower half of the
leaf. Nothing is written on it.}

\section{Leaf 66: Watercolors 1927 (clipping turned back)}

\sheet{16820}{66}

\hand{jo}{\emph{Water colors} \quad 1927 \quad Cape Elizabeth at Two Lights +
at Portland Head - while staying}

\begin{ledgertable}{Y{0.10}Y{0.06}Y{0.44}Y{0.09}Y{0.14}}{Delivered to Rehns & & & Amt. rec'd. & When rec'd.}
Oct. 21, 27 & 300 & Capt. Strout's House - Portland Head \quad 14 x 20 \quad \struck{300} Hartford Antheneum, Feb. 1928 & 200 & June 8, '28. \\
" & 300 & Foreshore - Two Lights \struck{surf}. Returned to Studio Nov. 12, 36. & & \\
" & 300 & Barn + Silo, Vt. (Hat on silo, barn on stilts) \quad 14 x 20 \quad Bgt. by Mr. L. G. Sheafer, Nov. 1927. & 200 & Mar. 9, 28. \\
" & 300 & Bill Latham's House (pine tree in foreground) house set back \quad 14 x 20 \quad Chas. Peppers, Nov. 3, 27 \quad Cape Elizabeth & 200 & Mar. 9, 28. \\
" & 300 & Two Lights Village (telegraph pole) \quad 14 x 20 \quad Returned + exchanged for Civil War Camp ground - wife did'nt like out house. \quad \struck{200 \quad Mar. 9, 28.} \quad - Dr. Lisle Nov. 11, 27 - \quad Henry Scherman Dec. 22, '28. (Book of Month) & 200 & Feb. 7, 29. \\
" & 300 & House of the Fog Horn II (from side in sunlight). \quad 14 x 20 \quad John T. Spaulding. Nov. 3, 27 \quad Cape Elizabeth \quad Boston Mus. 1948. & 200 & Mar. 9, 28. \\
Nov. 3, 27 & 300 & House of the Fog Horn I. grey day, rocks + cinders in fore-ground, flag pole - flag. \quad 14 x 20 \quad Cape Elizabeth \quad Mrs. John Osgord Blanchard. Nov. 11, 27. \quad Never rec'd money. Some mix up with th\ill{} & & \\
" & 300 & Hill + Houses. Capt. Berry's house foreground. Dark hill, \quad upton house, light house on top of hill. \quad 14 x 20 \quad John T. Spaulding - Nov. 3, 27 \quad Beauty \quad Cape Elizabeth \quad Boston Mus. 1948. & 200 & Mar. 9, 28. \\
" & 300 & Light at Two Lights. Blue sky, white light house back of upton House \quad House + out houses - movement \quad 14 x 20 \quad Miss Ruth Woodward, Nov. 3, 27. & 200 & Mar. 9, 28. \\
" & 300 & Coast Guard Station. Little station, very blue water of Cove. \quad very clean + swept look. \quad 14 x 20 \quad Leslie Sheafer, Nov. 3, 1927 \quad Two Lights & 200 & Mar. 9, 28 \\
" & 300 & Custom House - Portland - white pillars, windows, rain. \quad 14 x 20 \quad Huntington. Wadsworth Athenaeum via Huntingtons & 200 & Feb. 7, 29. \\
" & 300 & Libby House, Portland. Brownstone mansion + dead tree \quad trunk \quad 14 x 20 \quad Fog Museum, Dec. 12, 27 \quad Mr. Paul Sakes. Reprod. in Mod. Art particulars: book « Edward Hopper » 1933. & 200 & Mar. 9, 28. \\
" & 300 & Spurwink Church. grey day, churchyard in back. \quad 14 x 20 \quad Cape Elizabeth \quad Sold Ap. 23 1953 to Dr. Sheldon C. Sommers, Weston, Mass. \quad Harold C. Sommers \quad 400 - 1/3 & & Ap. 7, 1956 \\
" & 300 & Portland Head - Light. Blue silhouette - small + houses above \quad it base - gorgious blue day. \quad 14 x 20 \quad John T. Spaulding \quad Boston Mus. 1948. & 200 & Mar. 9, 28 \\
" & 300 & Rocky Pedestal - base of Portland Head Lighthouse on high \quad rocks, streaked up + down. Fog horn. \quad 14 x 20 \quad Nov. 12, 1936 to Frank Rehn as substitute for Tony's House purchased by him but sold + receipts went to E. H. & & \\
Nov. 12, 27 & 300 & Lobster Shack. Two Lights - in fog. \quad 14 x 20 \quad Cape Elizabeth \quad Mrs. John Osgord Blanchard - Nov. 18, 27 \quad 300 - 1/3 & 200 & Mar. 28, '30. \\
\end{ledgertable}

\hand{jo}{* Two Lights Village bgt. by Dr. Lyle Nov. 11, 27 but wife objected
to outhouse - so it was returned + exchanged for Civil War Camp Ground
(Rockland 1926)}

\marginal{Mod. Mus. cat - Edward Hopper 1933. \quad --- beside the House of the
Fog Horn I row}

\marginal{G. \quad --- above, between « Mr. L. » and « Sheafer »; Harry \quad
--- above, between « Chas. » and « Peppers »; ? \quad --- above « Ruth »;
church \quad --- above « graveyard »}

\marginal{Fr. Rehn \quad --- in pencil at the left of the Rocky Pedestal row}

\note{the asterisk that opens the note at the foot of the leaf answers a second
asterisk in the left margin beside the Foreshore row and a third at the right
of the Two Lights Village row; the note is hers and explains the struck
« 200 \quad Mar. 9, 28 » in that row. The doctor is « Dr. Lisle » in the row
and « Dr. Lyle » in the note, four inches apart on one leaf, and both stand.}

\note{« Osgord » is as written, confirming leaves 59 to 63. « Fog Museum » for
the Fogg and « Antheneum » / « Athenaeum » on two rows of this leaf are as
written. The words after « Some mix up with th » run into a browned strip at
the fore-edge and are readable in neither photograph of this leaf.}

\section{Leaf 67: Watercolors 1928, Gloucester}

\sheet{17040}{67}

\hand{jo}{\emph{Watercolors} 1928. Gloucester.}

\begin{ledgertable}{Y{0.10}Y{0.06}Y{0.44}Y{0.09}Y{0.14}}{ & & & & }
Oct 1928. & \$400 & Cape Anne Pasture. \quad 14 x 22. \quad Cows + rocks \quad longer shape. A beauty. \quad Stephen C. Clark - Oct. 1928. \quad 400 - 1/3 & 266 2/3 & June 5, 29. \\
 & & Adam's House - 16 x 25. Facade + roof at left. \quad \struck{\ill{}} porch overhanging with do-dads - Yellow hydrant, telegraph pole + wires, cloudy sky, st. + town beyond R.R. \quad Rehn - 400 \quad Gloucester \quad Presented to Frank + Peggy Rehn - Nov. 12, 1936. on the occasion of looking over the long list of sales made by them for E. H. & gift & \\
 & & Marty Welch's House - white, awnings + fence in sunlight. A beauty \quad 14 x 20 \quad Gloucester \quad John Clancey \quad fall 1928 \quad 300 - 1/3 & 200 & Mar. 29, 1930. \\
 & & House at Riverdale. \struck{\ill{}} red mansards + \quad 14 x 20 \quad Cape Anne \quad at L. fresh green tree swirling movement, rocky hill in back. \quad Whelan - \quad 300 - 1/3 & 200 & Mar. 28, 1930. \\
 & & My Roof - skylight + light over hall - grey glass \quad red tin rim. \quad 3 Washington Square, North, N.Y. City. \quad 14 x 20 \quad Dr. + Mrs. H. H. M. Lyle, surgeon, New York. \quad 300 - 1/3 & 200 & June 5, 29. \\
 & & Back St. Gloucester. Shacks, telegraph pole, bleak + grey in color. \quad 14 x 20 \quad Stephen C. Clark \quad 300 - 1/3 & 200 & June 5, 29. \\
\end{ledgertable}

\marginal{Reproduced in Whitney catalogue for Retrospective of E. H. 1950.
\quad --- beside the Adam's House row}

\marginal{full page reprod. in Mod. Mus. bk: Edward Hopper 1933. \quad ---
beneath the Marty Welch's House row}

\sketch{Back St. Gloucester \quad Stephen Clark}

\note{the Back St. Gloucester sketch is a pencil drawing on a small slip of
tinted paper pasted to the lower right of the leaf, in the gap between the
description and the money columns. It covers nothing.}

\note{the leaf heads no column. Two struck words open the Adam's House
description --- one written above the other and both crossed by the same
stroke --- and neither will resolve; the struck word in the House at Riverdale
description is a single word and will not resolve either. « Clancey » is as
written here against « Clancy » elsewhere in Book I.}

\section{Leaf 68: Watercolors 1928, Gloucester continued}

\sheet{18277}{68}

\hand{jo}{\emph{Watercolors} 1928 Gloucester - cont. \qquad 1928}

\begin{ledgertable}{Y{0.12}Y{0.50}Y{0.09}Y{0.14}}{ & & & When payment rec'd.}
Oct. 30, 28. & Sultry Day. Silver tones, houses, fresh dark green trees, \quad 14 x 20 \quad seen from high rock (that overlooks Anisquam River) \quad Leslie G. Sheafer - \quad 300 - 1/3 & 200 & Feb. 7, 29 \\
" & House on Middle St. (Gibson's) Porch + striped awning. \quad 14 x 20 \quad Door sized \quad (color of Marty Welch house) \quad Mrs. Sam Tucker - \quad 300 - 1/3 & 200 & Feb. 7, 29 \\
" & Box Factory. Covered bridge connecting 2 buildings, maroon color, st. in light below bridge. Box factory on R. not much in evidence. Auto middle distance. \quad 14 x 20 \quad Gloucester \quad Rep. Mod. Mus. cat. Edward Hopper 1933. \quad Mrs. John D. Rockefeller, Jr. - + given to Mod. Mus. that returned several watercolors of hers when they bought Corner Saloon - Bk. IV. \quad Bgt. by Mrs. Edward Marlin for 550. - 1/3 = \$366 2/3 \quad Feb. 14, 1957. & 200 & Feb. 7, 29 \\
" & Gloucester Roofs - lots of them - overlooked from st. above; st. of \quad Prospect St. looking toward Main St. Red mansard (white light on roof) + \quad \struck{Baptist bath}. on top. \quad 14 x 20 \quad Edward W. Root - about Oct. 31 - \quad 300 - & 200 & Feb. 7, 29. \\
Dec. 13, 28 & Circus Wagon - red wagon no. 16 - \quad Mrs. Swan. Jan. 21, 29 \quad 14 x 20 \quad Mrs. Swann formerly \quad Mrs. Paul Hammond \quad 300 - 1/3 & 200 & June 5, 29 \\
" & R.R. Gates. Black + white pointers, Catholic church spire, houses \quad 14 x 22? \quad Mrs. Samuel A. Tucker - 136 E. 79. Jan. 23, 29. \quad 400 - 1/3 & 266 2/3 & June 5, 29. \\
\end{ledgertable}

\marginal{Gloucester \quad --- above « Box Factory »; above, little \quad ---
above « st. above »; boards at \quad --- above « Red mansard »; Cupolo box
house \quad --- above the struck « Baptist bath »}

\note{« When payment rec'd. » is the only heading on this leaf, and « rec'd. »
is written above « payment » rather than after it. « Swan » and « Swann » stand
three words apart in the Circus Wagon row, the first in pencil and the second
in ink, and both are left as written; so is « Anisquam » for Annisquam.}

\section{Leaf 69: Watercolors, Charleston, S.C., April--May 1929}

\sheet{16592}{69}

\hand{jo}{\emph{Water colors.}\\
Charleston, S.C. \quad April, May 1929. \quad Three in gold frames + cream
mats.}

\marginal{1929 \quad Charleston, S.C. \quad Toppsfield, Mass., visiting Mrs.
Sam. Tucker. \quad Two Lights - last stay there.}

\begin{ledgertable}{Y{0.11}Y{0.50}Y{0.10}Y{0.14}}{ & & & }
May 25, 29 & Folly Beach (trees + sand - shaggy palms (trunks) + dead tree spreading arms. \quad 14 x 20 \quad Leslie Green Sheafer (May 1929) \quad 350 - 1/3 & 233 1/3 & Mar. 28, 30 \\
May 25, 29 & Negro Cabin. figure in dark doorway. Tall red chimney built outside. Fresh green trees to R. White cabin, black tar paper roof. \quad seated \quad 14 x 20 \quad Leslie Green Sheafer \quad 350 - 1/3 & 233 1/3 & Mar. 28, 30 \\
May 25, 29 & Charleston House - Tradd St. white shingles, red tiled roof, very brilliant white + sharp shadows - 2nd story, trees, trees front yard below. 3rd floor porch. Ash house. \quad 14 x 20 \quad Ash's House \quad to Rehn - fall 1929. \quad Wm. S. Paley, N.Y. \quad Wm. S.? \quad 29 Beekman Pl.? \quad 400 - 1/3 & 266 2/3. & May 26'' 1936 \\
May 25, 29. & Charleston Doorway - (+ gate posts, window with shutters at R. Brown + cream colored house - 1st story. \quad small \quad 14 x 20 \quad Mrs. Rockefeller - Oct. 29 \quad Mrs. John D. Rockefeller, Jr. \quad given by Mrs. R. to Mod. Mus. + returned by them in part payment for Corner Saloon - out of show of Early water colors \quad Smiths - \quad 400 - 1/3 = 266 2/3 & 266 2/3 & Nov. 10, 30. \quad June 16, 53. \\
May 25'', 29 & Baptistry of St. Johns. White marble rail + font, stair at R. brownish wall at back, tiled floor (grey + white squares) \quad 14 x 20 \quad (Church Interior after) \quad Mrs. Rockefeller. Oct. 29. \quad Mrs. John D. Rockefeller, Jr. + given to Mod. Mus \quad Mod. Mus. returns it in part payment for Corner Saloon. \quad Goldstone - 600 - 1/3 Com. = 400 \quad July 8, 1957. & 266 2/3 & Nov. 10, '30. \\
\end{ledgertable}

\sketch{}

\note{four of Edward Hopper's ink record drawings stand on this leaf, ruled
into boxes at the left of the money columns --- one each for Folly Beach, Negro
Cabin, Charleston House and Charleston Doorway, and a fifth for the Baptistry.
None carries any writing. An ink asterisk stands in the right margin beside the
Charleston Doorway row; its sense is not established.}

\note{the leaf heads no column. « Toppsfield » for Topsfield is as written; so
is the second « Wm. S.? » in pencil under the first. The Charleston Doorway row
carries two dates in its last column, « Nov. 10, 30. » in ink and
« June 16, 53. » in blue, the second belonging to the return to the Museum of
Modern Art described in the same row.}

\section{Sheet of drawings laid over leaf 69}

\sheet{18337}{}

\note{a separate sheet, narrower than the leaf, lying over leaf 69 --- the
Whitney's descriptor says adhered to it. Eleven of Edward Hopper's ink
thumbnail drawings stand on it, in three groups, with Jo Hopper's sale lines
beneath them. The heading and the first date column of leaf 69 show above and
to the left of it, and its right-hand money column shows at the fore-edge;
both are transcribed at sheet 16592 and are not repeated here.}

\hand{jo}{\emph{Drawings}}

\sketch{}

\begin{ledgertable}{Y{0.44}Y{0.22}Y{0.17}}{ & & }
Miss Russell \quad Feb. 2''. & 100 - 1/3. \quad 66 2/3 & Rec'd. June 5, 29 \\
Gift to Forbes + Nan Watson & & \\
Withdrawn. & & \\
Nan Watson - paid & & \\
\add{t}o Rehn Gallery - Jan. 14, 29. \quad Hopper Ex. of Oils, Water Colors + Drawings & & Ap. 26'' 1936. \\
\end{ledgertable}

\hand{jo}{Charleston, S.C. \quad April, May 1929.}

\begin{ledgertable}{Y{0.44}Y{0.22}Y{0.17}}{ & & }
Drawing: Negro Houses - May 25'', 1929. Framed + delivered to Gallery. \quad Conté crayon - \quad Bought by Tom Colt - \quad Colt - of the (Richmond Mus.) & \$73 2/3 & Mar. 28, 1930. \\
Drawing: Gloucester Houses at \quad Bgt. by \quad McCurrach - & 100 - 1/3 = 66 2/3 - & Rec'd June 15, 44. \\
Circus Tents - Gift to Mrs. Sam. A. Tucker - & & June 1929 at Toppsfield \\
Landscape drawing - to Santa Barbara Museum. & 100 - 1/3 - \$66 2/3. & June 27, 49. \\
Group of Cats - (Perkins Dos Passos' « Siamese ») bgt. by \quad Haswell \quad (have High Noon oil) & \$100 - 1/3 = 66 2/3 & Dec. 29, 51. \\
House, Gloucester \quad Interior of Theatre \quad Corcoran \quad 100 \quad 60 & 160 - 1/3 = \$106 2/3 & Dec. 29, 51. \\
Drawing. Trees, Maine - Mrs. James H. Beal & 100 - 1/3 = 66 2/3, & July 5, 1955. \\
\end{ledgertable}

\note{« Withdrawn. » stands under the standing nude of the first group and
« Nan Watson - paid » under the leftmost of the six figure drawings of the
second group; neither is a title. The « t » of « to Rehn Gallery » runs off the
left edge of this sheet and is supplied. « House, Gloucester » and « Interior
of Theatre » are braced together as one sale, the two prices 100 and 60 bracing
to the 160 beside them. The figure for the Negro Houses drawing is « \$73 2/3 »
as written, and is not reconciled with the third commission the line records.}

\section{Leaf 70: Water Colors, Summer 1929}

\sheet{16348}{70}

\hand{jo}{\emph{Water Colors} - Summer 1929. \quad Two Lights, Cape Elizabeth,
Me.}

\begin{ledgertable}{Y{0.10}Y{0.06}Y{0.44}Y{0.09}Y{0.14}}{ & & & & }
Oct. '29 & 400 & Coast Guard Boat I. \quad 14 x 20 \quad A beauty. Blue sky, white boat, blue water, sand. \quad Mr. Huntington. Rob't. W. \quad of Hartford, Conn. \quad Mr. Huntington, Hartford. Spring 1930. \quad Rep. Mod. Mus. cat. E. H. 1933 & 266 2/3 & Oct. \quad Nov. 10, '30. \\
 & 400 & Coast Guard Boat II. \quad 14 x 20 \quad L. Temple & 500 - 1/3 = 333 1/3. & Ap. 7, 1956. \quad frame. \\
 & 400 & Coast Guard Cove \quad 14 x 20. \quad A beauty - pale blue water, cliff in shadow - light on little lookout shack. white boat \quad Mr. + Mrs. Geo. H. Davis \quad 400 - 1/3 \quad 5'' Ave. \quad Small collection of Henri, Lux, \struck{\ill{}} & 266 2/3. & Mar. 5, 31. \\
 & 400 & Boat and Cliff \quad 14 x 20 \quad Coast guard boat in profile against cliff - boat small \quad Barney - fall 1929 - \quad 400 - 1/3 \quad Bought at Auction by \uncertain{Frank Cahn} - 1944 - \uncertain{p. 325} & 266 2/3 - & Mar. 28, 1930 \\
 & 400 & Rocky Cove \quad 14 x 20 \quad near \uncertain{Maginello} \quad much in \struck{\ill{}} \quad silhouette. \quad Returned to Studio Nov. 12, 1936 to use frame on new picture. \quad Rolf Dessauer \quad 400 - 1/3 & 266 2/3, & July 5, 1955. \\
 & 400 & The Dory \quad 14 x 20 \quad Dory with \struck{sails} on roughish grey blue water - rocks foreground - yellow. \quad Rector of Princeton \quad Dr. Bryan 1929. \quad Williams B. \quad Dr. Bryan Jr. (D.D.) of Princeton - minister of & 233 1/3 & Nov. 10, '30 \\
\end{ledgertable}

\marginal{silhouette \quad --- above the struck word in the Rocky Cove row;
deep - yellowish \quad --- see leaf 71}

\note{the leaf heads no column. « Maginello » will not settle and neither will
the pencil page reference at the end of the Boat and Cliff row, offered here as
« p. 325 »; « Frank Cahn » is doubtful in its second word only. The word struck
after « Henri, Lux, » is one of her own scribbled cancellations and will not
recover at any magnification. Four of Edward Hopper's ink record drawings stand
in the money column of this leaf, for Coast Guard Boat I, Coast Guard Boat II,
Coast Guard Cove and Boat and Cliff, and two more below; « Coast Guard Cove »
and « Boat + Cliff » are written under two of them in her hand.}

\note{the last row runs off the foot of the leaf: « minister of » is the last
thing the photograph shows and the sentence is not finished on this sheet.}

\sketch{Coast Guard Cove}

\sketch{Boat + Cliff}

\section{Leaf 71: Summer 1929, Cape Elizabeth and Essex}

\sheet{17836}{71}

\hand{jo}{Summer 1929.\\
Cape Elizabeth + Trips to Essex \quad House at Essex \quad Barn at " \quad
+ Pemaquid - Pemaquid Light}

\begin{ledgertable}{Y{0.10}Y{0.06}Y{0.44}Y{0.09}Y{0.14}}{ & & & & }
Oct. '29 & 400 & House at Essex \quad 14 x 20 \quad green fields, spick white house, stone fence \quad spick green tree in front of house. \quad Mrs. Murray Crane (Modern Mus.) Mar. 1930 \quad Mrs. Murray Crane - 400 - 1/3 \quad N.Y. & 266 2/3 - & Mar. 5, 31. \\
 & 400 & Methodist Church \quad 14 x 20 \quad (Church Yard. Cape Eliz. \quad From side rear of graveyard - 1 grave. Square tower of church \quad Mrs. Shepherd - fall 1929 \quad Mr. + Mrs. Shepherd, John S. \quad New York. & 266 2/3 & Nov. 10, '30 \\
 & 400 & House of the Fog Horn III - 14 x 20. \quad Very bright colored \quad Dr. Lisle - Nov.? 1929. \quad Dr. H. H. M. Lyle N.Y.C. - surgeon. \quad - Hopper Fitch of Yale - now owns. \uncertain{Met} of Whitney. June 1956. & 266 2/3 & Nov. 10, 30. \quad 8'' \\
 & 500 & Light House Village \quad 16 x 25 \quad Palish blue sky, long fluffy clouds \quad brilliant grass in foreground \quad Cleveland Museum - Mar. 1930. \quad \ill{} & 266 2/3 & Mar. 5, 31. \\
\ill{} & \ill{} & \ill{} & 333 1/3 & Nov. 10, 30 \quad 8'' \\
\ill{} & \ill{} & \ill{} \quad y day (white day) \quad 750 - 1/3 = 500. & & Dec. 29, 51. \\
\ill{} & \ill{} & Dr. Harold Sommers, Newton Center (June 30, 50) & & \ill{} \\
\end{ledgertable}

\marginal{Farm \quad --- above « House at Essex »; stretch of \quad --- above
« long fluffy clouds »; deep - yellowish \quad --- above « in foreground »;
church \quad --- above « graveyard »}

\clipping{a halftone of a watercolour --- a white barn with a round wooden
silo, a car in the yard, houses beyond --- on a wide card mount, hinged at its
left edge and lying over the lower left of the leaf, captioned « FARM AT
ESSEX » over « Edward Hopper ».}

\note{the clipping covers the whole of two entries and the descriptions that
went with them. The Whitney photographed this leaf a second time with the
clipping turned back --- sheet 16955, which it calls « Page 71 - Verso » and
which is the recto again --- and that sheet is the first of the next batch of
twelve. The three \ill{} rows above carry what stands clear of the clipping's
right edge and nothing else; each \ill{} there means « under paper », the
mount reaching the leaf's left margin so that the date and price columns of
those rows are covered too. The last \ill{} of all is a short run of ink marks
at the very foot of the leaf, cut by the leaf's own edge. The rows themselves
are to be read at 16955, and the four entries already given above must not be
transcribed a second time.}

\note{the leaf heads no column. « Farm » written above « House at Essex » is
the title the clipping's caption prints, so the correction is hers and predates
the cutting; the entry is transcribed as the ink gives it. The « 333 » of the
fifth row is written over an earlier figure. The « 8'' » that follows two of
the dates is as written and its sense is not established. The word before « of
Whitney » will not settle, and « Met » is offered rather than read.}

\keywords{medium:watercolours, medium:drawings, place:Rockland,
place:Gloucester, place:Cape Elizabeth, place:Two Lights, place:Portland Head,
place:Charleston, place:Essex, place:Cape Anne, dealer:Frank K. M. Rehn,
person:Mrs. John Osgord Blanchard, person:John Taylor Spaulding,
person:Stephen C. Clark, person:Mrs. John D. Rockefeller Jr.,
person:Edward W. Root, person:Leslie Green Sheafer, person:Frank
Crowninshield, person:Adolph Lewisohn, person:Mrs. Murray Crane,
person:William S. Paley, person:Robert Wheelwright, person:Alfredo Valente,
person:A. Percy Saunders, person:Hubert Goldstone, collection:Boston Museum
of Fine Arts, collection:Museum of Modern Art, collection:Wadsworth
Atheneum, collection:Cleveland Museum, collection:Fogg Art Museum,
collection:Hamilton College, collection:Santa Barbara Museum,
collection:Corcoran Gallery, collection:New Britain Museum, feature:dealer
commissions, feature:resales, feature:exchanges, feature:gifts,
feature:record sketches, feature:pasted clippings}

\end{document}
